\documentclass[10pt,a4paper]{amsart}
\usepackage[margin=19mm]{geometry}
\usepackage{amsmath,amssymb,mathtools}
\usepackage[scr=rsfs,cal=euler]{mathalfa}
\usepackage{xfrac}
\usepackage[unicode,hidelinks,bookmarks=false]{hyperref}
\newcommand{\ie}{i.e.\ }
\newcommand{\cf}{cf.\ }
\newcommand{\resp}{respectively\ }
\newcommand{\Subsection}{\subsection}
\providecommand{\nobreakdash}{\nobreak\hbox{-}}
\setlength{\emergencystretch}{2em}
\multlinegap0pt
\usepackage{amssymb}
\usepackage{tikz}
\newtheorem{theorem}{Theorem}
\newcommand{\round}{{\rho}}
\title{Burning Steiner Triple Systems}
\author{Andrea C. Burgess, Peter H. Danziger, Caleb W. Jones, Trent G. Marbach, David A. Pike}
\date{Source: arXiv:2608.25627v1 (CC BY 4.0)}
\begin{document}
\maketitle

\noindent\emph{Source and licence.} Burning Steiner Triple Systems. Version arXiv:2608.25627v1 (CC BY 4.0), distributed by arXiv under the Creative Commons Attribution 4.0 International licence, http://creativecommons.org/licenses/by/4.0/, as recorded in the arXiv licence metadata for this exact version and shown on the abstract page https://arxiv.org/abs/2608.25627v1. Full licence notice: \texttt{licenses/arxiv-2608.25627-CC-BY-4.0.txt}.\par\smallskip

\noindent\textit{Editorial scope.} Counted unit: the article's theorem h(r) (source0.tex lines 344--349). The article's complete original proof of it is delivered with it (source0.tex lines 350--380, one \texttt{proof} environment of 31 lines). No mathematics is written, repaired or re-derived: the statement and the proof are the article's own lines, in the article's own notation, and no retained line is substituted or re-flowed.  Retained uncounted as the source-local context the proof needs: lines 273--290.  Nothing was omitted from any retained span, so the package carries no omission record.  No retained line contains a citation, so the wrapper carries no bibliography and no bibliography entry.  No retained line contains a figure environment, an image-inclusion command or drawing code, so no image is delivered and the package declares no asset dependency.  Editorial formatting only, in addition: an AMS \texttt{amsart} preamble that restates the article's own declarations and macros used by the retained lines, namely the environments \texttt{theorem} on separate counters, together with the standard packages those lines need.  The retained environments print in this document as Theorem 1 in order of appearance, whereas the article prints the same items with the numbering of its own full document.  The complete native source of the article as distributed in the arXiv e-print for this exact version is bundled unchanged as \texttt{sources/208588/source0.tex}.
\begin{table}[ht]
\centering
\begin{tabular}{|c|c|}
\hline
$\round$ & $h(\round)$   \\ \hline 
1 & 1          \\ \hline  
2 & 2          \\ \hline  
3 & 4          \\ \hline  
4 & 8          \\ \hline  
5 & 25         \\ \hline  
6 & 266        \\ \hline  
7 & 34\,732      \\ \hline 
8 & 603\,069\,104  \\ \hline  
9 & 181\,846\,170\,592\,008\,673 \\ \hline 
\end{tabular}
\caption{The first nine values for the function $h$.}
\label{h_g_table}
\end{table}

\begin{theorem}
If $\round \geq 9$, $\alpha = 1.07925087759784$, and 
$\beta = 1.07925087759785 = \alpha+10^{-14}$, then
    $$2 \alpha^{2^{\round}}+\frac{5}{2} \leq h(\round) \leq 2\beta^{2^{\round}}.$$
\label{Size of h(r)}
\end{theorem}

\begin{proof}
We proceed by induction on $\round$.  
For the base case for the upper bound, consider $\round=9$.  It is straightforward to verify that $\lfloor 2 \beta^{2^\round}\rfloor =181846170592709201$, which exceeds $h(9)$; see Table~\ref{h_g_table}. 

Now, recalling that $h(\round)=\binom{h(\round-1)}{2}-2h(\round-1)+3\round-2$, we have the inductive step, as 
\begin{align*}
h(\round) &= \binom{h(\round-1)}{2} -2h(\round-1)+3\round-2 \\
     &= h(\round-1)\left( \frac{h(\round-1)-1}{2} - 2 \right) +3\round-2 \\
     &\leq 2\beta^{2^{\round-1}} \left( \beta^{2^{\round-1}} - \frac{5}{2}\right)+3\round-2\\
     &= 2\beta^{2^{\round-1}\cdot 2} - 5 \beta^{2^{\round-1}} +3\round-2\\
     & \leq 2\beta^{2^{\round-1}\cdot 2}\\
     & = 2\beta^{2^{\round}},
\end{align*}
where the last inequality holds since $5 \beta^{2^{\round-1}} \geq 3\round-2$ for $\round \geq 5$. This completes the inductive step, and the proof for the upper bound. 

For the lower bound, we consider as the base case $\round=8$. It follows from calculation that $\lceil 2 \alpha^{2^\round}+5/2\rceil = 603069104$. 
This is equal to $h(8)$; see Table~\ref{h_g_table}. 


For the inductive step, we have 
\begin{align*}
h(\round) &=\binom{h(\round-1)}{2} -2h(\round-1)+3\round-2 \\
     &= h(\round-1)\left( \frac{h(\round-1)-1}{2} - 2 \right) +3\round-2 \\
     &\geq \left( 2 \alpha^{2^{\round-1}} +\frac{5}{2} \right)\left( \alpha^{2^{\round-1}}-\frac{5}{4}\right)+3\round-2\\
     &= 2 \alpha^{2^{\round-1}\cdot 2} - \frac{25}{8} +3\round-2\\
     & \geq 2 \alpha^{2^{\round}}- \frac{5}{2}+3\round-3+\frac{3}{8}\\
     & \geq 2 \alpha^{2^{\round}}-\frac{5}{2},
\end{align*}
where we have $3\round\geq 3-\frac{3}{8}$ for $\round \geq 1$. 
This completes the inductive step, and the proof for the lower bound.
\end{proof}

\end{document}
